%
% nullstellen.tex -- 
%
% (c) 2021 Prof Dr Andreas Müller, OST Ostschweizer Fachhochschule
%
\documentclass[tikz]{standalone}
\usepackage{amsmath}
\usepackage{times}
\usepackage{txfonts}
\usepackage{pgfplots}
\usepackage{csvsimple}
\usetikzlibrary{arrows,intersections,math,calc}
\definecolor{darkred}{rgb}{0.8,0,0}
\definecolor{darkgreen}{rgb}{0,0.6,0}
\begin{document}
\def\skala{1}
\begin{tikzpicture}[>=latex,thick,scale=\skala,x=0.9cm,y=0.12cm,font=\small]
		
	% ------------------------------------------------------------
	% Achsen
	% ------------------------------------------------------------
	\draw[->] (-8.8,0) -- (4.2,0) node[right] {$\operatorname{Re}(s)$};
	\draw[->] (0,-34) -- (0,38) node[above] {$\operatorname{Im}(s)$};
	
	% ------------------------------------------------------------
	% Kritischer Streifen 0 < Re(s) < 1
	% ------------------------------------------------------------
	\fill[gray!36,opacity=0.5] (0,-30) rectangle (1,34);
	\draw[thick] (0,-30) -- (0,34);
	\draw[thick] (1,-30) -- (1,34);
	
	% kritische Gerade Re(s)=1/2
	\draw[dashed,thick] (0.5,-30) -- (0.5,34);
	
	% ------------------------------------------------------------
	% Markierungen auf der reellen Achse
	% ------------------------------------------------------------
	\foreach \x in {-8,-6,-4,-2,0,1}
	{
		\draw (\x,0.8) -- (\x,-0.8);
	}
	
	\draw (0.5,0.6) -- (0.5,-0.6);
	
	\node[below] at (-8,-0.8) {$-8$};
	\node[below] at (-6,-0.8) {$-6$};
	\node[below] at (-4,-0.8) {$-4$};
	\node[below] at (-2,-0.8) {$-2$};
	
	\node[below left]  at (0,-0.8) {$0$};
	\node[below left]       at (0.5,-0.6) {$\frac12$};
	\node[below left] at (1,-0.8) {$1$};
	
	% ------------------------------------------------------------
	% triviale Nullstellen
	% ------------------------------------------------------------
	\foreach \x in {-8,-6,-4,-2}
	{
		\fill[blue] (\x,0) circle (2.1pt);
	}
	
	%\node[align=center] at (-6.1,-12)
	%{
	%	triviale\\
	%	Nullstellen
	%};
	
	%\draw[->,green] (-5.95,-9.2) -- (-2.02,-0.5);
	%\draw[->,green] (-5.95,-9.2) -- (-4.02,-0.5);
	%\draw[->,green] (-5.95,-9.2) -- (-6.02,-0.5);
	%\draw[->,green] (-5.95,-9.2) -- (-8.02,-0.5);
	% ------------------------------------------------------------
	% nichttriviale Nullstellen
	% echte erste paar Ordinaten, alle auf Re(s)=1/2
	% ------------------------------------------------------------
	\foreach \y in {14.1347,21.0220,25.0109,-14.1347,-21.0220,-25.0109}
	{
		\fill[green!70!black] (0.5,\y) circle (2.2pt);
	}
	
	% ------------------------------------------------------------
	% Beschriftung kritischer Streifen
	% ------------------------------------------------------------
	\node[align=left] at (2.1,18)
	{
		kritischer\\
		Streifen
	};
	
	% kritische Gerade
	\node[align=left] at (3.5,7)
	{
		kritische Gerade\\
		$\operatorname{Re}(s)=\frac12$
	};
	
	\draw[->] (1.95,6.5) -- (0.56,6.5);
	
	% nichttriviale Nullstellen
	%\node[align=right] at (-1.5,23)
	%{
	%	nichttriviale\\
	%	Nullstellen
	%};
	
	%\draw[->,blue] (-0.55,21.5) -- (0.44,-14.1);
	%\draw[->,blue] (-0.55,21.5) -- (0.44,-21.0);
	%\draw[->,blue] (-0.55,21.5) -- (0.44,-25.0);
	%\draw[->,blue] (-0.55,21.5) -- (0.44,25.0);
	%\draw[->,blue] (-0.55,21.5) -- (0.44,21.0);
	%\draw[->,blue] (-0.55,21.5) -- (0.44,14.1);
	% ------------------------------------------------------------
	% Pol bei s=1
	% ------------------------------------------------------------
	\draw[red,very thick] (0.90,-1.0) -- (1.10,1.0);
	\draw[red,very thick] (0.90,1.0) -- (1.10,-1.0);
	
	\node[red,anchor=west] at (1.45,-4.0) {Pol bei $s=1$};
	\draw[->,red] (1.38,-3.2) -- (1.05,-0.5);
	
\end{tikzpicture}
\end{document}

